\chapter{The Hypothesis AST}
\label{ch:hypothesis}

Every hunt is driven by a parsed \code{Hypothesis} value. The AST
is deliberately small --- a flat vector of steps with per-step
predicate lists, plus optional aggregation and time-window clauses
attached to the whole chain.

\section{Types}

\begin{lstlisting}[style=rust]
pub struct Hypothesis {
    pub name:           String,
    pub steps:          Vec<HypothesisStep>,
    pub k_simplicity:   usize,            // default 1
    pub aggregations:   Vec<AggClause>,   // HAVING  clauses
    pub window_seconds: Option<u64>,      // WITHIN  clause
}

pub struct HypothesisStep {
    pub origin_type:     EntityType,
    pub relation_type:   RelationType,
    pub dest_type:       EntityType,
    pub edge_predicates: Vec<Predicate>,  // on edge metadata
    pub dest_predicates: Vec<Predicate>,  // on destination-vertex metadata
    pub origin_binding:  Option<String>,  // `(User u)` for ref in HAVING
    pub dest_binding:    Option<String>,
}

pub enum Predicate {
    Eq      { key: String, value: String },
    Substr  { key: String, pattern: String },
    Regex   { key: String, pattern: String },
    In      { key: String, values: Vec<String> },
}

pub struct AggClause {
    pub op:     AggOp,        // Count, CountDistinct, Min, Max, Sum, Avg
    pub column: AggColumn,    // bound name or step index
    pub cmp:    AggComparison, // >=, <=, >, <, ==
    pub value:  f64,
}
\end{lstlisting}

\section{Invariants}

\begin{invariant}[Step count $\geq 1$]
A hypothesis must have at least one step; the DSL parser rejects
the empty input with \code{ApiError::InvalidInput}. This keeps the
matcher from having to treat the zero-edge case as a special
branch.
\end{invariant}

\begin{invariant}[$k$-simplicity positive]
$k_{\text{simplicity}} \geq 1$. Zero would admit only the empty
path and violates the sense of ``max visits per vertex.''
\end{invariant}

\begin{invariant}[Binding uniqueness]
If \code{origin\_binding} or \code{dest\_binding} is \code{Some(b)}
on any step, $b$ does not repeat on any other binding position
inside the hypothesis. The parser enforces this at parse time so
the matcher can look bindings up by name during aggregation
without disambiguation.
\end{invariant}

\section{Static compatibility check}

Before running a hunt, the engine can cross-reference the
hypothesis against the graph's relation schema to see if any step
is structurally impossible on the loaded dataset --- e.g., a
\code{User -[Auth]-> Host} step against a purely cloud-identity
graph where no \code{Host} entities exist. This is
\textsc{Check-Catalog-Compatibility} (Chapter~\ref{ch:catalog-yaml}),
which returns a \code{CatalogEntryWithStatus} flagging the missing
entity or relation types per step.

\begin{algorithm}
\caption{\textsc{Check-Step-Compatibility}}\label{alg:check-step}
\KwIn{schema $\Sigma$ (multiset of $(\text{source}, \text{rel},
\text{target})$ triples), step $s$}
\KwOut{\code{bool} compatible, set of missing axes}
\For{each $(\text{src}, r, \text{dst}) \in \Sigma$}{
    \If{$(\text{src} = s.\text{origin\_type} \vee s.\text{origin\_type} = *)$
       $\wedge (r = s.\text{relation\_type} \vee s.\text{relation\_type} = *)$
       $\wedge (\text{dst} = s.\text{dest\_type} \vee s.\text{dest\_type} = *)$}{
        \Return $(\mathrm{true}, \emptyset)$\;
    }
}
$\text{missing} \gets $ set of axes absent from any triple in $\Sigma$ \;
\Return $(\mathrm{false}, \text{missing})$\;
\end{algorithm}

\begin{complexity}
$\Ord{|\Sigma|}$ per step, $\Ord{|\Sigma| \cdot \ell}$ per
hypothesis, where $\ell$ is the step count. $|\Sigma|$ is bounded by
$|\mathcal{T}_V|^2 \cdot |\mathcal{T}_E|$ in practice (there are
only a few dozen entity/relation types), so the check is
effectively constant.
\end{complexity}

\begin{inthecode}
\filepath{graph\_hunter\_core/src/hypothesis.rs} --- \code{Hypothesis},
\code{HypothesisStep}, \code{Predicate}.\\
\filepath{graph\_hunter\_core/src/catalog/mod.rs:105} ---
\code{check\_catalog\_compatibility}.
\end{inthecode}
